% ============================================================================
% Appendix: experiment summary and development trace.
% Sourced from report/analysis/appendix_experiment_summary/ (audit, tables,
% paragraphs). All public/private Macro F1 values trace to the Kaggle API
% (docs/submissions_archive/submission_manifest.csv); no score is inferred from
% a training log. Requires \usepackage{tabularx} and \usepackage{booktabs}.
% ============================================================================

\appendix
\section{Experiment Summary and Development Trace}
\label{app:experiment_summary}

This appendix collects every major direction we explored for PlantCLEF 2026, what
each one tested, how deeply we pursued it, its best leaderboard scores, and what we
took away from it. It is intended as an honest development trace rather than a
controlled ablation study. Over the project we made more than six hundred Kaggle
submissions; the tables below summarise them by experiment or by inference direction
rather than listing every run. No score is inferred
from a training log. Where a private score could not be recovered we write
\texttt{[nf]} (not found), and directions that never produced a submission (data
construction and diagnostic analyses) are marked \texttt{[ns]}.

\paragraph{How to read ``exploration depth''.}
We label each direction with one of five depths.
\emph{Deep} means many runs or sweeps, follow-up analysis, or a component of the
final system (for example the i002 model, the aggregation choice, and the headline
ensemble). \emph{Moderate} means several meaningful runs but not an exhaustive sweep
(for example the early baselines). \emph{Shallow} means one or a few quick trials.
\emph{Diagnostic only} means an analysis used to understand behaviour rather than to
chase a score. \emph{Dropped} means a direction we explored and then abandoned after
weak or near-neutral results. These labels describe effort and role, not statistical
significance.

% ---------- Table 1: model / training lineage ----------
\begin{table}[t]
\centering
\scriptsize
\setlength{\tabcolsep}{4pt}
\renewcommand{\arraystretch}{1.15}
\caption{Model and training experiments. Public/private are the
best Kaggle Macro F1 for that experiment group; for the model rows the private value
is the best private of the group. Depth: Deep / Moderate / Shallow / Diagnostic / Dropped. Scores are not a
controlled ablation; see the caveats in Section~\ref{sec:appendix_caveats}.}
\label{tab:appendix_main_models}
\begin{tabularx}{\linewidth}{@{}l l c c c c X@{}}
\toprule
\textbf{Exp.} & \textbf{Direction / main change} & \textbf{Pub.} & \textbf{Priv.} & \textbf{Depth} & \textbf{Fwd?} & \textbf{Main lesson} \\
\midrule
001/002 & BioCLIP zero-shot tiling (no training) & $\sim$0.12 & $\sim$0.12 & Mod. & concept & Biologically useful but far too weak zero-shot; needs a trained head. \\
003 & BioCLIP 2.5 linear probe (frozen) & 0.242 & 0.281 & Mod. & No & Doubles zero-shot but plateaus; frozen backbone cannot bridge the domain shift. \\
004 & BioCLIP few-shot support bank & 0.307 & 0.296 & Mod. & No & Prototypes beat the probe but stay well short of fine-tuning. \\
005 & DINOv3 SSL + fine-tune & 0.140 & 0.142 & Shal. & No & DINOv3 SSL-adapt scored far below BioCLIP fine-tunes. \\
006 & BioCLIP 2.5 partial fine-tune (linear head) & 0.222 & 0.272 & Mod. & partial & Unfreezing raised val but not public; needed a richer head + better inference. \\
007 & BioCLIP zero-shot + SAM vegetation mask & 0.140 & 0.142 & Drop. & No & SAM-weighted aggregation did not improve on plain zero-shot ($\sim$70~h inference). \\
008 & DINOv3+PlantNet fine-tune (no local folder) & 0.375 & 0.351 & Mod. & No & Strongest non-BioCLIP line, still below 010/i002. \\
009 & BioCLIP 2.5 + PlantNet data & 0.208 & 0.248 & Shal. & No & PlantNet data worse than PlantCLEF 2024 single-plant. \\
015 & PlantCLEF24 + iNaturalist mix (no folder) & 0.380 & 0.342 & Shal. & No & Higher val, lower public: a validation/leaderboard mismatch case. \\
010 & BioCLIP 2.5 end-to-end multitask, \emph{shared} MLP, 5 heads & 0.391 & 0.389 & Deep & Yes & End-to-end multitask jumps to $\sim$0.39; softmax-mean + last-4-blocks are load-bearing. \\
011 & 010 + SimSiam SSL on pseudo-quadrats & 0.377 & 0.356 & Mod. & No & SSL hurt the frozen-head stage and never beat 010. \\
014 & Unfrozen-blocks sweep (no folder) & 0.375 & 0.343 & Diag. & informs & Confirms an interior optimum in unfrozen-block count. \\
i001 & Data download / manifest construction & [ns] & [ns] & Deep & Yes & Larger clean genus/family-filled manifest enabled the i002 jump. \\
\textbf{i002} & \textbf{BioCLIP 2.5, \emph{per-head} MLPs, more data, no cap} & \textbf{0.412} & \textbf{0.406} & Deep & \textbf{Yes} & Final model: more clean data + per-head MLPs + softmax-mean give the project best. \\
i003 & i002 + $148$k GBIF images for $<\!100$-image species, $500$/species cap & 0.400 & 0.402 & Deep & No & Capping/flattening hurt public; more clean data beat a balanced smaller set. \\
\bottomrule
\end{tabularx}
\end{table}

% ---------- Table 2: inference, aggregation, ensembles, pivots, diagnostics ----------
\begin{table}[t]
\centering
\scriptsize
\setlength{\tabcolsep}{4pt}
\renewcommand{\arraystretch}{1.15}
\caption{Inference, aggregation, ensemble, pivot, and diagnostic directions, all built on
the i002 (or 010) checkpoint. Public/private are the best Kaggle Macro F1 for the named
configuration. The 224+336 ensemble is the reported headline; the best private overall is
the i002 logit-adjusted run. Diagnostics produced no submission ([ns]). This is a
development trace, not a controlled ablation; see Section~\ref{sec:appendix_caveats}.}
\label{tab:appendix_inference}
\begin{tabularx}{\linewidth}{@{}l l c c c c X@{}}
\toprule
\textbf{Item} & \textbf{Direction / main change} & \textbf{Pub.} & \textbf{Priv.} & \textbf{Depth} & \textbf{Fwd?} & \textbf{Main lesson} \\
\midrule
Agg. & softmax-mean vs max vs mean (pooling) & 0.412 & 0.406 & Deep & Yes & softmax-mean beats max ($+0.037$) and avoids mean dilution; the production aggregator. \\
mean & plain logit-mean (010 head-only only) & 0.258 & 0.320 & Drop. & No & Dilutes a species strong in 1--2 tiles across background; dropped before i002. \\
LA & logit adjustment $\tau$ sweep & 0.406 & \textbf{0.406} & Deep & Yes & Mild $\tau{=}0.25$ lowers public slightly but gives the best private; $\tau{\geq}0.5$ harmful. \\
Sel. & selection rule (probT / top-$k$ / relT / gap) & 0.406 & 0.406 & Deep & Yes & Probability thresholding ($T{=}0.03$) at $\tau{=}0.25$ beats top-$k$/gap/relT (no-LA peak is 0.412 at $T{=}0.05$). \\
\textbf{Ens.} & \textbf{224+336 resolution ensemble (headline)} & \textbf{0.418} & 0.403 & Deep & \textbf{Yes} & Multi-resolution average of one checkpoint: best public overall. \\
Ens.2 & i002 + 010 ensemble & 0.389 & 0.405 & Shal. & No & Hurt public, near-best private: vivid public/private disagreement. \\
Piv.1 & cRT triple-backbone classifier & 0.382 & [nf] & Drop. & No & A different classifier diverged into wrong answers, not productive diversity. \\
Piv.2 & genus/family co-occurrence rerank & 0.412 & 0.399 & Drop. & No & Priors from the model's own posterior add no information; near-neutral. \\
Piv.3 & seasonal phenology prior (novelty) & 0.413 & 0.388 & Deep & novelty & Date is orthogonal in theory but not statistically; visual model already absorbs season. \\
Diag.1 & head-architecture audit (010 vs i002) & [ns] & [ns] & Diag. & report & No clean ablation isolates head topology or taxonomy supervision. \\
Diag.2 & head-feature CKA analysis & [ns] & [ns] & Diag. & report & Per-head MLPs are distinct, not redundant; family MLP most taxonomy-structured. \\
Diag.3 & within-backbone saturation gates & [ns] & [ns] & Diag. & informs & All within-i002 variants are essentially one model; headroom needs orthogonal signal. \\
\bottomrule
\end{tabularx}
\end{table}

\subsection{Baselines and zero-shot experiments}
We first established that the BioCLIP family of foundation models carried useful
biological signal, then measured how far that signal went without task-specific
training. Zero-shot tile classification (experiments 001 and 002), in which test
tiles are matched against species text prompts, scored only around $0.12$ public
Macro F1 at best, and the SAM vegetation-weighted variant (007) did not improve on
plain zero-shot aggregation while costing roughly seventy hours of inference. A
frozen-backbone linear probe (003) roughly doubled the zero-shot score to $0.242$
public but then plateaued, and a non-parametric few-shot support bank (004) reached
$0.307$. The lesson from this group is consistent: BioCLIP embeddings are a strong
starting point, but neither zero-shot matching nor a frozen backbone is enough to
bridge the shift from single-plant training images to multi-species quadrats.

\subsection{Supervised fine-tuning experiments}
The decisive gains came from partially fine-tuning BioCLIP 2.5 ViT-H/14. A
single-head partial fine-tune (006) improved validation accuracy but not public
Macro F1, which signalled that we needed both a richer head and a better inference
pipeline. Experiment 010 introduced the end-to-end multitask design: a shared
multi-layer perceptron feeding species, genus, family, order, and class heads, with
fixed auxiliary loss weights. It reached $0.391$ public and made two choices that
carried through the rest of the project, namely softmax-mean tile aggregation and
unfreezing only the last few transformer blocks. Experiment 011 added SimSiam
self-supervised pretraining on pseudo-quadrats, but this hurt the frozen-head stage
and never beat 010. The strongest model, i002, kept the 010 recipe but trained on the
larger and cleaner i001 manifest (about $2.65$M images), switched to independent
per-head multi-layer perceptrons for species, genus, and family, and trained for ten
epochs. Its best single-configuration public score was $0.412$ and it produced the
best private score of the whole project, $0.40600$, under a lightly logit-adjusted
configuration. Experiment i003 then capped the data at five hundred images per
species and added extra GBIF images; this flatter distribution \emph{lowered} public
Macro F1 to $0.400$, so it was not carried forward. We emphasise that the gap between
010 and i002 mixes several changes (head topology, data, head count, epochs, and batch
size) and is therefore not attributable to any single factor.

\subsection{Inference and aggregation experiments}
On a fixed checkpoint we explored how to turn sixteen per-tile predictions into one
species set. Among the three pooling rules, softmax-mean (averaging per-tile softmax
probabilities) was the clear choice: on the i002 four-block checkpoint it beat
per-species max by about $0.037$ public Macro F1, with the same ordering on the
private split. Plain logit-space mean was explored only on the weaker 010 head-only
checkpoint, where it was consistently the worst of the three (best public $0.258$),
because averaging across sixteen cells dilutes a species that is strong in only one or
two tiles; we therefore dropped it before i002. For species selection, an adaptive
probability threshold ($T \approx 0.03$ to $0.05$) dominated fixed top-$k$, relative
threshold, and gap rules. A mild class-prior logit adjustment ($\tau{=}0.25$) traded a
little public score for the best private score. Finally, averaging the same checkpoint
at two inference resolutions ($224$ and $336$ pixels) produced our headline public
result of $0.418$. A cross-checkpoint ensemble with 010 instead hurt public but
reached near-best private ($0.405$).

\subsection{Data experiments}
Experiment i001 rebuilt the training corpus by re-downloading PlantCLEF 2024,
deduplicating an iNaturalist research-grade dump, and adding GBIF images for
under-represented species, then back-filling genus and family for every row. The
resulting larger, cleaner manifest is what the winning i002 model trained on, with no
per-species cap applied. Two data variations were tested and rejected on the
leaderboard: mixing in iNaturalist data (group 015) raised validation but lowered
public Macro F1, and capping at five hundred images per species (i003) also lowered
public score. The extra GBIF images did not move any species above the hundred-image
threshold, so the long tail remained long. The takeaway is that more clean data helped,
but artificially balancing or expanding the distribution did not.

\subsection{Diagnostic analyses}
Three analyses were run to understand the system rather than to chase a score. A head
architecture audit documented that 010 uses a shared trunk while i002 uses independent
per-head multi-layer perceptrons, and established that no clean ablation isolates the
head topology or the taxonomy supervision, so neither can be claimed as the cause of
i002's gain. A representation analysis (linear CKA and silhouette on validation
features) showed that the i002 per-head perceptrons are genuinely distinct rather than
redundant, with the family head most taxonomy-structured, while all heads retain the
BioCLIP geometry. A within-backbone saturation diagnostic showed that variants of i002
(different resolutions and unfreeze depths) agree on roughly eighty to ninety percent
of top-1 predictions, so the remaining internal diversity is below the leaderboard
noise floor; this motivated the search for orthogonal signals in the three pivots.

\subsection{Final system lineage}
The final system emerged along a clear path, which we describe as a sequence of choices
that coincided with our best results rather than as a chain of proven causes. Zero-shot
and frozen-backbone baselines established that BioCLIP was useful but insufficient.
End-to-end multitask fine-tuning in experiment 010 produced the first strong model and
fixed two recurring ingredients, softmax-mean aggregation and last-few-block
unfreezing. The i001 data rebuild and the move to per-head perceptrons in i002 (trained
on more clean data with no cap) gave the strongest checkpoint. On top of that
checkpoint, the aggregation study selected softmax-mean over max and dropped plain mean,
the selection study chose adaptive probability thresholding, light logit adjustment
($\tau{=}0.25$) gave the best private score, and a $224{+}336$ resolution ensemble gave
the best public score of $0.41826$. The three pivots (a triple-backbone classifier, a
genus/family rerank, and a seasonal phenology prior) were all reranking attempts on this
anchor; none improved it, and we report the phenology prior as a novelty contribution
on its own terms.

\subsection{Important caveats for interpreting the appendix}
\label{sec:appendix_caveats}
\begin{itemize}
  \item This is a development trace, not a controlled ablation table. Most rows change
  several factors at once (data, architecture, schedule, and inference), so a score
  difference between two rows cannot be attributed to any single change.
  \item In particular, the 010-to-i002 improvement is confounded by the data rebuild,
  the head redesign, the head count, the epoch count, and the batch size. We did not
  run a controlled shared-versus-per-head comparison, nor a species-only run that
  removes the taxonomy heads, so we report these as design choices, not causes.
  \item Public and private leaderboards disagree. Our best public configuration
  ($0.41826$) is not our best private one ($0.40600$), and the i002+010 ensemble is
  near-best on private while weak on public. We report both columns and avoid implying
  that the best public run is best overall.
  \item Held-out single-plant validation accuracy correlates only weakly with
  multi-species quadrat Macro F1 (about $r{\approx}0.4$ across our strongest runs). The
  four-block checkpoint beats the eight-block one on the leaderboard despite lower
  validation accuracy, so validation numbers should not be over-read.
  \item Some directions are shallow (one or a few submissions) or have no recoverable
  local code, especially the manifest-only groups 008, 009, 014, and 015. Their setup
  is inferred from submission messages and should not be over-interpreted.
  \item The diagnostic analyses are descriptive. They characterise the model and its
  representations but do not establish causal effects on the score.
  \item Missing private scores are marked \texttt{[nf]} and data-only or diagnostic rows
  are marked \texttt{[ns]}. Several near-neutral pivot deltas (about $-0.005$ to
  $-0.006$) are small, about two to three times the leaderboard noise floor we estimated at
  roughly $0.002$, and should be read as ``no useful effect'' rather than as a measured decrease.
\end{itemize}
